\section{开始之前：把图像问题读成数字问题}
\subsection{本书的学习顺序}
一张灰度照片可以看作一张数字表。表中的每个数表示一个小位置的亮暗；把这些位置按原来的排列显示出来，便成为图像。图像处理的大部分入门操作，都可以从这个事实出发理解：改变一个数，可以改变一个像素的亮暗；把周围几个数平均，可以减小局部起伏；比较相邻位置的数，可以发现边缘。

学习过程分为四步。第一步理解图像是怎样记录位置与亮暗的，以及人眼怎样观察这些变化。第二步研究对单个灰度值作变换，并统计整张图的灰度分布。第三步让相邻像素参与计算，理解平滑、去噪和锐化。第四步把图像理解为不同快慢变化的组合，进入频率域滤波。最后一种表示并没有创造另一张与原图无关的图，而是把同一份信息换一种方式描述。

本书沿用附件中的题号，例如E1、D4，使原题与详解可以相互定位。但知识讲解与对应例题放在一起，不要求先读完全部概念再集中做题。题目中的条件、计算结果和适用范围必须一起理解。原稿中的“必考”“高频考点”等标签不是本书对考试命题的保证。

\subsection{符号只是名称：先明确输入与输出}
字母常用来给一个数或一组数起名字。例如，用$r$表示处理前的灰度，用$s$表示处理后的灰度。若规则是“在原数上加20”，就写成
\[
s=r+20.
\]
当$r=30$时，$s=30+20=50$；当$r=100$时，$s=120$。字母没有增加新的计算，只是避免为每个可能的输入重新写一遍同样的话。

函数（Function）是从输入得到输出的规则。把上述规则叫作$T$，可以写成$s=T(r)$，其中$T(r)=r+20$。这里的括号表示“将$r$交给规则$T$处理”，不是相乘。图像处理中常见的$f(x,y)$也遵循这个习惯：$x,y$指明位置，$f(x,y)$给出该位置的灰度值。它不表示$f$乘$x$再乘$y$。

下标（Subscript）用于区分同类对象。$r_0,r_1,r_2$表示第0、第1、第2个灰度值；$n_k$表示第$k$类的数量。下标的作用像座位编号，它本身不是乘法。上标有时表示幂，例如$r^2=r\times r$，有时是标记；具体含义必须看定义。

\subsection{分数、平均与比例}
分数线表示除法。$\frac{a}{b}$表示$a\div b$，前提是$b\ne0$。如果5个像素的灰度分别是10、20、20、30、40，它们的总和是120，平均值为
\[
\frac{10+20+20+30+40}{5}=\frac{120}{5}=24.
\]
平均值回答的是“如果把总量均匀分给这5个位置，每个位置得到多少”。它不要求24恰好是原来5个数之一。

比例回答的是“某一部分占整体多少”。若100个像素中有20个灰度为0，灰度0的比例为$20/100=0.2=20\%$。后文把这种比例称为相对频数；在图像的经验统计模型中，也把它当作该灰度的概率（Probability）。这里所需的概率概念首先就是计数与除法，不必先学习复杂的随机过程。

\subsection{求和符号不是一种新运算}
希腊大写字母$\Sigma$表示把一串数相加。例如
\[
\sum_{i=1}^{4}a_i=a_1+a_2+a_3+a_4.
\]
符号下面的$i=1$给出起点，上面的4给出终点。$i$是依次取1、2、3、4的编号。若$(a_1,a_2,a_3,a_4)=(2,3,5,6)$，则结果为16。

二维图像既有行又有列，因此有时出现两层求和。例如一个$2\times2$小区域的4个数为
\[
\begin{bmatrix}1&2\\3&4\end{bmatrix}.
\]
先把每一行相加，得到3和7，再把行的结果相加，得到10。双重求和符号只是把这个“逐行、逐列遍历”的过程缩写起来。后文的卷积与相关公式，都是先给出实际小网格算例，再写这种缩写形式。

\subsection{平方、平方根与指数}
平方是同一个数乘自己：$3^2=3\times3=9$。平方根（Square Root）是反过来寻找非负的那个数：$\sqrt9=3$。例如直角三角形的两条直角边分别为3和4，斜边长度是$\sqrt{3^2+4^2}=5$。频谱图中某点离中心有多远，也用同一种距离计算。

指数（Exponent）表示乘法的重复或其连续扩展。$2^3=2\times2\times2=8$，$2^{-1}=1/2$，$r^{1/2}=\sqrt r$。伽马变换中的$r^\gamma$就是这种幂运算；判断它如何改变亮暗，可以先代入$r=0.25$：$0.25^{1/2}=0.5$，而$0.25^2=0.0625$。在0至1的灰度尺度上，前者变亮，后者变暗。

自然指数函数写作$\exp(x)=e^x$，其中$e\approx2.71828$。对数（Logarithm）是指数运算的逆运算：在$a>0$、$a\ne1$、$c>0$的条件下，若$a^b=c$，就有$\log_a c=b$。自然对数记为$\ln c$，以$e$为底。特别重要的一条规则是
\[
\ln(ab)=\ln a+\ln b,\qquad a>0,\ b>0.
\]
这解释了后面同态滤波为什么先取对数：它希望把“照明乘反射”变成“两项相加”。在使用这条公式前，必须确保要取对数的数为正。

\subsection{单位是计算的一部分}
一个数后面的单位说明它在量什么。毫米描述长度，像素描述离散采样点，毫米/像素描述每个采样点对应的实际长度。若10毫米宽的纸面由20个像素记录，则
\[
\frac{10\ \mathrm{mm}}{20\ \mathrm{pixel}}=0.5\ \mathrm{mm/pixel}.
\]
长度单位必须统一后再相加或比较。常用换算为
\[
1\ \mathrm{m}=1000\ \mathrm{mm},\qquad
1\ \mathrm{mm}=1000\ \mu\mathrm{m}.
\]
例如$4.52\,\mu\mathrm{m}=0.00452\,\mathrm{mm}$。如果把4.52直接当作毫米代入，相机视场会被放大1000倍；这种错误不是算法问题，而是量纲不一致。

分辨率计算常需要“向上取整”。若计算出至少333.33个像素，实际不能制造三分之一个像素，也不能取333而低于下限，所以至少取334。向上取整记为$\lceil333.33\rceil=334$。普通四舍五入与向上取整不是同一种操作，必须按题意选择。

\subsection{示意图、精确计算与现实测量}
示意图用于突出关系，例如两个像素是否相邻、一条变换曲线位于恒等线的哪一侧。它不一定按真实尺寸比例绘制。精确计算则必须明确输入数据、公式和假设。现实测量还受到噪声、模糊、几何误差等因素影响，因此“计算得到足够像素”不自动等于“实际一定能识别目标”。

后文会区分这三个层次。需要理解机制时使用小整数或小网格；求数值答案时给出具体单位和逐步代入；涉及设备选择时再说明理想模型的边界。这样可以先建立清楚的基本关系，再逐渐加入真实场景的复杂因素。
